%=====================================================================
\section{The environment: formulation 3 as a Markov decision process}
\label{sec:mdp}

The environment (\code{rl2\_env.py}, identical to RL1's) is a batched PyTorch simulator: it plays many episodes side by side, for many instances or many rollouts of one instance, with every quantity a tensor of shape (batch, \dots). This section defines it mathematically.

\subsection{Resources, state and initial state}

Let $G$ be the set of guns ($|G|=24$ in the 100-EV files) with station $f(g)$, and $M$ the set of vans ($|M|=5$). The resources are $K=G\cup M$. The state before decision $t$ is
\[
s_t=\Big(P_t,\;D_t,\;(a_g)_{g\in G},\;(\ell_m,a_m,b_m,\kappa_m)_{m\in M}\Big),
\]
where $P_t$ is the set of placed EVs, $D_t$ the dropped EVs, $a_g$ the minute gun $g$ becomes free, and for van $m$: $\ell_m$ its last node (0 = depot), $a_m$ the minute it is free there, $b_m$ its battery left and $\kappa_m$ the EVs it has served. Initially $P_0=D_0=\emptyset$, $a_g=0$, $\ell_m=0$, $a_m=0$, $b_m=B$, $\kappa_m=0$.

\subsection{Actions, timing and masks}

An action is a pair $(i,k)$ with $i\notin P_t\cup D_t$ and $k\in K$. Its start time, wait and finish follow Assumption~\ref{as:timing}:
\begin{align*}
\text{gun } g:&\quad S_i=\max\!\big(r_i+t_{if(g)},\,a_g\big),\quad q^\Fcs_i=S_i-r_i-t_{if(g)},\quad C_i=S_i+\tau^\Fcs_i;\\
\text{van } m:&\quad S_i=\max\!\big(r_i,\,a_m+t_{\ell_m i}\big),\quad q^\Mcs_i=S_i-r_i,\quad C_i=S_i+\tau^\Mcs_i .
\end{align*}
The pair is \textbf{feasible} (unmasked) if
\begin{align*}
\text{gun } g:&\quad d_{if(g)}\le R_i,\qquad q^\Fcs_i\le W_i,\qquad C_i\le T,\qquad \nexists\,g'<g:\ f(g')=f(g),\ a_{g'}=a_g;\\
\text{van } m:&\quad \kappa_m<Q,\qquad b_m\ \ge\ \frac{E_i}{\eta}+\beta\big(d_{\ell_m i}+d_{i0}\big),\qquad C_i+t_{i0}\le T,\\
&\quad \kappa_m>0\ \text{ or } m \text{ is the first van with } \kappa=0 .
\end{align*}
Each condition is one constraint of formulation 3 checked at the moment of placing: reachable station (the set $G_i$), queue tolerance (eq.\ 5), deadline (eq.\ 21), seats (eq.\ 15), battery with the way home (eq.\ 16), van back by $T$ (eq.\ 19); the last condition of each block is the symmetry reduction (Assumption~\ref{as:sym}). The van battery condition reserves the energy to drive home from $i$; because of the triangle inequality, later EVs on the route can only add to that reserve, never invalidate it.

An EV is \textbf{selectable} if it has at least one feasible pair. Before every decision, EVs that are not placed and have no feasible pair are moved to $D_t$ (Proposition~\ref{prop:dead} shows they could never be placed later). The episode ends when $P_t\cup D_t=I$.

\subsection{Transition}

After action $(i,k)$: $P_{t+1}=P_t\cup\{i\}$, and
\begin{align*}
\text{gun } g:&\quad a_g\leftarrow C_i;\\
\text{van } m:&\quad b_m\leftarrow b_m-\frac{E_i}{\eta}-\beta\,d_{\ell_m i},\quad \ell_m\leftarrow i,\quad a_m\leftarrow C_i,\quad \kappa_m\leftarrow\kappa_m+1 .
\end{align*}
The simulator is deterministic: the same action sequence always gives the same plan. RL2's active search relies on this to replay the incumbent (Section~\ref{sec:eas}).

\subsection{Reward}
\label{sec:reward}

The reward is $-\Delta Z$ with
\begin{align*}
\Delta Z(i,g)=&\ \mu_i\big(c^{\mathrm{EV}}d_{if(g)}+p^\Fcs E_i\big)+\pi^\Fcs E_i+c^{\mathrm D}_i t_{if(g)}+c^{\mathrm Q}_i q^\Fcs_i+c^{\mathrm C}_i\tau^\Fcs_i,\\
\Delta Z(i,m)=&\ \mu_i\,p^\Mcs E_i+\pi^\Mcs\frac{E_i}{\eta}+c^{\mathrm V}_i\big(q^\Mcs_i+\tau^\Mcs_i\big)+c^\Mcs\big(d_{\ell_m i}+d_{i0}-d_{\ell_m 0}\big)+c^\Fcs\,[\kappa_m=0].
\end{align*}
The van term charges the detour: the van used to end its route by driving $\ell_m\to0$; now it drives $\ell_m\to i\to0$.

\begin{proposition}[The rewards add up to the objective]\label{prop:reward}
If an episode places every EV, the sum of its $\Delta Z$ equals $Z$ of eq.~\eqref{eq:Z} for the resulting plan.
\end{proposition}
\begin{proof}
$Z_1$, $Z_4$, $Z_5$ and $Z_6$ are sums of per-EV terms, and each EV's terms appear in $\Delta Z$ exactly once, when it is placed. $Z_3$: the dispatch cost $c^\Fcs$ is charged when a van serves its first EV ($\kappa_m=0$), so once per used van. $Z_2$: for a van with route $0\to i_1\to\cdots\to i_k\to0$ the van terms are
\[
c^\Mcs\big[(d_{0i_1}+d_{i_10}-d_{00})+(d_{i_1i_2}+d_{i_20}-d_{i_10})+\dots+(d_{i_{k-1}i_k}+d_{i_k0}-d_{i_{k-1}0})\big],
\]
which telescopes (with $d_{00}=0$) to $c^\Mcs\big(d_{0i_1}+d_{i_1i_2}+\dots+d_{i_{k-1}i_k}+d_{i_k0}\big)$, the length of the route.
\end{proof}

\subsection{Three properties that justify the design}

\begin{proposition}[Earliest start is optimal]\label{prop:timing}
Fix the resource of every EV and the order of the EVs on every gun and van. If some start times $S$ satisfy all constraints of formulation 3, then the earliest-start times $S^{*}$ of the simulator satisfy them too, and $Z(S^{*})\le Z(S)$.
\end{proposition}
\begin{proof}
With assignment and sequences fixed, the constraints on $S$ are (a) lower bounds: $S_i\ge r_i+t_{if}$ (FCS) or $S_i\ge r_i$ and $S_i\ge t_{0i}$ for the first EV of a van; $S_j\ge S_i+\tau^\Fcs_i$ for consecutive EVs on a gun; $S_j\ge S_i+\tau^\Mcs_i+t_{ij}$ for consecutive EVs of a van; and (b) upper bounds: $S_i-r_i-t_{if}\le W_i$, $S_i+\tau_i\le T$, $S_i+\tau^\Mcs_i+t_{i0}\le T$ for the last EV of a van. $S^{*}$ sets every $S_i$ to the largest of its lower bounds, going along each sequence. By induction along the sequences, any $S$ satisfying (a) has $S\ge S^{*}$ componentwise. Every upper bound in (b) has the form $S_i\le\text{const}$, so if $S$ satisfies it, so does $S^{*}\le S$. Finally, $S$ enters $Z$ only through $c^{\mathrm Q}_iq^\Fcs_i$ and $c^{\mathrm V}_iq^\Mcs_i$ with $q=S_i-\text{const}$ and non-negative coefficients, so $Z$ is non-decreasing in every $S_i$ and $Z(S^{*})\le Z(S)$.
\end{proof}

\begin{proposition}[No plan is out of reach]\label{prop:complete}
For every feasible plan $\pi$ of formulation 3 there is a sequence of unmasked actions whose resulting plan has the same assignment and sequences as $\pi$ (up to exchanging identical guns of one station or identical vans) and a cost no larger than $Z(\pi)$.
\end{proposition}
\begin{proof}
Order the EVs by their start time in $\pi$ (ties by position in their sequence) and append them in that order to their resource in $\pi$. Each resource's sequence is then rebuilt in its own order, and when EV $i$ is appended its predecessor is the last EV of that resource. By induction, the simulator's start times are not later than $\pi$'s (Proposition~\ref{prop:timing} applied to the prefix), so the queue, deadline and return conditions hold. The van seat count never exceeds $\pi$'s. For the battery: in $\pi$ the energy of the whole route fits in $B$; the energy still needed after reaching $i$ is $E_i/\eta$ plus the energy of the rest of the route, and by the triangle inequality the rest of the route home is at least $d_{i0}$ long, so the mask condition $b_m\ge E_i/\eta+\beta(d_{\ell_m i}+d_{i0})$ holds. If the symmetry rule hides the gun or van used in $\pi$, an identical one is offered; exchanging the two labels for the rest of $\pi$ changes nothing. The cost is at most $Z(\pi)$ by Proposition~\ref{prop:timing}.
\end{proof}

Proposition~\ref{prop:complete} is why the action space is a \emph{pair} (EV, resource): the optimal plan is always reachable. Had the EVs been processed in a fixed order (e.g.\ by request time) and only the resource been chosen, plans in which a later request is served before an earlier one on the same resource would be out of reach.

\begin{proposition}[A dead EV stays dead]\label{prop:dead}
If an unplaced EV $i$ has no feasible pair in state $s_t$, it has none in any later state.
\end{proposition}
\begin{proof}
Guns: $a_g$ never decreases, so $S_i$, $q^\Fcs_i$ and $C_i$ never decrease and the range condition never changes; a gun hidden by the symmetry rule has the same $a_g$ as the offered one. Used vans: after van $m$ serves some $j$, $\kappa_m$ grows; its start for $i$ is $\max(r_i,C_j+t_{ji})\ge\max(r_i,a_m+t_{\ell_m j}+t_{ji})\ge\max(r_i,a_m+t_{\ell_m i})$ by the triangle inequality, so it is not earlier; and the battery shortfall $E_i/\eta+\beta(d_{ji}+d_{i0})-b'_m$ with $b'_m=b_m-E_j/\eta-\beta d_{\ell_m j}$ exceeds the old shortfall by $E_j/\eta+\beta(d_{\ell_m j}+d_{ji}-d_{\ell_m i})>0$. Unused vans are identical and unchanged. So every condition that failed still fails.
\end{proof}

\subsection{One step of the simulator}

\begin{algorithm}[H]
\caption{One construction step (all rows of the batch at once)}
\DontPrintSemicolon
compute for every unplaced EV $i$ and resource $k$: $S_i$, $q_i$, $C_i$, feasibility, $\Delta Z(i,k)$ \tcp*{\code{Env.options}}
move EVs with no feasible pair to $D$ \tcp*{\code{Env.drop\_dead}}
\If{all EVs are placed or dropped}{\Return}
policy chooses EV $i$ among selectable EVs, then resource $k$ among feasible resources of $i$\;
cost $\leftarrow$ cost $+\Delta Z(i,k)$; record resource, $S_i$, $C_i$, step number of $i$\;
update $a_g$ (gun) or $(\ell_m,a_m,b_m,\kappa_m)$ (van) \tcp*{\code{Env.step}}
\end{algorithm}

\subsection{From an episode to a plan}

\code{plan\_of} turns one finished episode into the plan format of the project: the resource of every EV (``FCS$f$-gun$g$'' or ``van$v$'', vans in use numbered $1,2,\dots$ as in the MILP), start and completion times, and each van's route $0\to\dots\to0$. Dropped EVs are completed as in Assumption~\ref{as:complete}. The plan then goes through \code{results\_io.evaluate}, which calls \code{check\_plan\_f3} and writes \code{runs.csv} and \code{plan.csv}. The simulator's $Z$ is compared with the checker's and a warning is written if they differ by more than $10^{-6}$ relative.
